\SourceLeaf{73}
Hinc sponte sequitur

\[
S - \frac{dT}{dx} = \left(\frac{\partial W}{\partial r}\right),
\]
\[
R - \frac{dS}{dx} + \frac{d^2T}{dx^2} = \left(\frac{\partial W}{\partial q}\right),
\]
\[
Q - \frac{dR}{dx} + \frac{d^2S}{dx^2} - \frac{d^3T}{dx^3} = \left(\frac{\partial W}{\partial p}\right),
\]
\[
P - \frac{dQ}{dx} + \frac{d^2R}{dx^2} - \frac{d^3S}{dx^3} + \frac{d^4T}{dx^4} = \left(\frac{\partial W}{\partial y}\right),
\]
\[
N - \frac{dP}{dx} + \frac{d^2Q}{dx^2} - \text{etc.} = 0.
\]

Ex iisdem principiis theorema inversum nullo negotio demonstrare licet.

\noindent (Sept.\ 1800.)

\par\smallskip\hrulefill\smallskip\par

\Needspace{10\baselineskip}\noindent\hspace*{\parindent}Den Beweis des Satzes, dass die Bedingungsgleichung

\[
\left[N - \frac{dP}{dx} + \frac{d^2Q}{dx^2} - \frac{d^3R}{dx^3} + .\,.\,. = 0\right]
\]

\noindent statthat, wenn $Vdx$ ein vollständiges Differential ist, habe ich in meinem Exemplar von \textsc{Euler}s Integralrechnung [Band III, Seite 518] gegeben. Der Beweis des [umgekehrten] Satzes beruht auf folgenden Momenten.

Da der Ausdruck $N - \frac{dP}{dx} + \frac{ddQ}{dx^2} - \text{etc.}$ auf eine ganz bestimmte Art aus der Natur der Function $V$ abgeleitet wird, so bezeichnen wir ihn durch $fV$, und eben so $P - \frac{dQ}{dx} + \frac{ddR}{dx^2} - \text{etc.}$ durch $f'V$ u.\,s.\,w.

Man setze $\int f'V  .  dy = W^0$, so integrirt, dass bloss $y$ als veränderlich betrachtet wird, und das vollständige Differential $dW^0 = V^0dx$. Dann wird\HistoricalNote{*)}{*)}{Diese beiden Schlüsse hängen mit dem Beweise a.\,a.\,O.\ in \textsc{Euler}'s Integralrechnung zusammen, wonach, $W^0$ statt $W$ gesetzt, $f'V^0 = \dfrac{\partial W^0}{\partial y}$ also $= f'V$ wird, und $fV^0 = 0$ also $= fV$.}

\[
\left.\begin{array}{l}
f'V^0 = f'V \\
fV^0 = fV
\end{array}\right\}
\text{ und daher }
\begin{array}{l}
f'(V - V^0) = 0, \\
f(V - V^0) = 0.
\end{array}
\]

\noindent Da nun

\[
\frac{\partial(V - V^0)}{\partial y} = f(V - V^0) + \frac{df'(V - V^0)}{dx}
\]

\noindent [gilt,. so ist auch $\dfrac{\partial(V - V^0)}{\partial y} = 0$] und $V - V^0$ enthält kein $y$.

